\subsection{Task 3.2. Nhật ký dự án hàng tuần (Weekly Project Log)}
\label{sec:t3-weekly}

Nhật ký dưới đây mở rộng Bảng \ref{tab:i1-weekly} cho toàn bộ thời gian thực hiện. Giờ công kế hoạch lấy từ Bảng \ref{tab:ledger} (Mục 2.11). Công việc thực tế được đối chiếu với lịch sử commit và các biên bản kiểm thử trong repository. Do nhóm không ghi timesheet đều đặn từ đầu dự án, giờ công thực tế (đánh dấu $^{*}$) là \textbf{ước tính hồi cứu} theo phương pháp ở Mục \ref{sec:t3-effort-method}, không phải số đo trực tiếp. Tuần được tính từ thứ Hai đến Chủ nhật; tuần T1 bắt đầu ngày 26/10/2026.

\begingroup\footnotesize
\renewcommand{\arraystretch}{1.18}
\begin{longtable}{@{}p{.11\linewidth}p{.20\linewidth}p{.25\linewidth}p{.19\linewidth}p{.19\linewidth}@{}}
\caption{Weekly Project Log T1--T6}\label{tab:t3-weekly}\\
\toprule
\textbf{Tuần / giờ} & \textbf{Kế hoạch (V1.0)} & \textbf{Thực tế (bằng chứng)} & \textbf{Vấn đề / rủi ro} & \textbf{Hành động quản lý}\\\midrule\endfirsthead
\toprule
\textbf{Tuần / giờ} & \textbf{Kế hoạch (V1.0)} & \textbf{Thực tế (bằng chứng)} & \textbf{Vấn đề / rủi ro} & \textbf{Hành động quản lý}\\\midrule\endhead
\midrule\multicolumn{5}{r}{\textit{Còn tiếp}}\\\endfoot
\bottomrule\endlastfoot
\textbf{T1}\newline 26/10--01/11\newline KH: 63,5h\newline TT$^{*}$: 56,5h &
G02--G04 lập và duyệt baseline; A01 thiết kế; A02.1 test case I1; A03 môi trường (29--30/10); bắt đầu A05.2 (30/10). &
Baseline V1.0 hoàn tất; ERD, đặc tả API, mockup; bản đặc tả 24 use case (31/10). A03 chưa xong trong tuần. &
A03 trễ: Docker Desktop, socket Java trên Windows, múi giờ JDBC (AR3). &
Giữ nhánh thiết kế/test chạy song song; dời A05.2 sang sau khi có môi trường.\\
\textbf{T2}\newline 02--08/11\newline KH: 64,0h\newline TT$^{*}$: 79,5h &
A05.1, A05.2 lập trình I1; A04.1 dữ liệu thử; A06 kiểm thử I1; RA02 thử Lightsail; RA03 thử email; RA04 thử QR. &
Smoke test local 03/11 (\texttt{b57b1aa}); tài khoản/RBAC 03/11 (\texttt{2ed54e3}); API sự kiện 04/11 (\texttt{bf9cc50}); UI I1 04--06/11; regression 07/11: 208 passed/1 skipped; PR \#1 merge 08/11. &
A03 trễ 4 ngày; RA02 chưa làm; lỗi BUG-P2-05; QA phát hiện sức chứa Unlimited và lọc thời gian lệch đặc tả. &
Chốt phạm vi 03--04/11: bỏ UC-24 (CR-01), sức chứa nguyên dương (CR-03); ưu tiên sửa BUG-P2-05.\\
\textbf{T3}\newline 09--15/11\newline KH: 61,1h\newline TT$^{*}$: 69,7h &
CP1 (09/11); A07.1, A07.2; RA01 thử tải đồng thời; RA06 làm rõ yêu cầu; BL1 dữ liệu thử I2/I3; A18.1. &
PR \#2 merge 12/11; tái cấu trúc backend phân tầng (\texttt{4d56706}); backend I2 hoàn chỉnh cùng ngày (\texttt{43f8c4c}), 29 unit test và HTTP test đạt; thử gửi 1 email thật qua Brevo. I1 được hợp nhất vào \texttt{main} ngày 12/11 (PR \#2), trễ 3 ngày so với mốc CP1 (09/11). &
09--11/11 không có commit, I2 bắt đầu muộn; tái cấu trúc là việc ngoài kế hoạch. &
Nhóm tự chốt quy tắc FIFO, email ngày 12/11 (chưa có xác nhận của khách hàng); gộp A07.1--A07.3 vào một đợt để bù thời gian.\\
\textbf{T4}\newline 16--22/11\newline KH: 60,2h\newline TT$^{*}$: 65,0h &
A07.3 email; A08 kiểm thử I2 (CP2); A09.1, A09.2 điểm danh; A10.1 báo cáo; A15.1; A19.1. &
Backend I3 18/11 (\texttt{1b99579}), 43 test đạt; frontend I2 merge 18/11 (PR \#3); trang điểm danh Organizer 18--19/11; báo cáo và CSV 22/11 (\texttt{98f6024}). I2 được hợp nhất ngày 18/11 (PR \#3), sớm 1 ngày so với mốc CP2 (19/11). &
I2/I3 nằm trên nhánh tính năng, \texttt{main} vẫn ở I1; test email phụ thuộc thời gian chạy (flaky). &
Chốt giờ mở QR và bước xác nhận 18/11; sửa test harness email, không sửa logic.\\
\textbf{T5}\newline 23--29/11\newline KH: 53,4h\newline TT$^{*}$: 35,5h\newline (đến 27/11) &
Hoàn tất A09.1, A10.2; A11 tích hợp, triển khai (CP3); A12.1 integration test; A12.2 đo chỉ tiêu; A14.1 User Guide; RA05. &
Trang quét QR 23/11 (\texttt{26efcbb}); rà soát toàn hệ thống và merge vào \texttt{main} 25/11 (\texttt{727f379}); deploy Vercel/Render/Neon 25/11; kiểm tra cloud 27/11: health 200. Nhận CR-05 (rút ngắn 1 tuần). A12.1 đang thực hiện. &
Đổi nền tảng so với Lightsail; Maven 429 khi build; request đầu sau khi backend ngủ mất $\approx$ 50 giây; mất toàn bộ buffer lịch do CR-05. &
Lập CR-04; dùng mirror Maven + retry; làm nóng backend trước demo (NR1). Phân tích CR-05, chọn PA3: chạy song song A12 và báo cáo, dồn T6 vào 27--29/11.\\
\textbf{T6}\newline 30/11--06/12\newline KH: 29,4h\newline (dồn vào T5 theo CR-05) &
A12.2, A15.2, A16, A18.2, A19.2, BL2; demo A17 (03/12); nộp báo cáo 06/12. &
Các hoạt động T6 được thực hiện trong 27--29/11 theo kế hoạch CR-05 (Bảng \ref{tab:t3-cr05-plan}); ước tính 37,0 giờ theo Bảng \ref{tab:t3-realloc}. &
Không còn buffer; review của khách hàng chưa thực hiện (NR7). &
Demo trình bày change log để khách hàng quyết định trực tiếp các CR.\\
\midrule
\textbf{Tổng} & \multicolumn{4}{l}{Kế hoạch 331,6 giờ công (T1--T5: 302,2). Thực tế$^{*}$ T1--T5 đến 27/11: 306,2 giờ công.}\\
\end{longtable}
\endgroup

Điểm đáng chú ý nhất của nhật ký là \textbf{nhịp làm việc không đều}: các tuần có khoảng trống không có commit (09--11/11), sau đó khối lượng lớn được hoàn thành dồn vào một hai ngày (12/11, 18/11, 25/11). Ngày hoàn thành trên git vì vậy phản ánh thời gian trôi (elapsed time), không phản ánh trực tiếp công sức (effort). T2 vượt kế hoạch 15,5 giờ do môi trường trễ và giao diện I1 phải làm lại; T5 thấp vì A12.2, A14.1 chưa bắt đầu tại ngày báo cáo.
